\chapter{Configuration, Commands, and Glossary}

Keep this appendix near your terminal. It gathers names and commands that are easy to forget after you understand the larger idea. Copy a command only after reading the short note around it, especially when hardware is connected.

\section{Environment variables}

\begin{longtable}{>{\ttfamily\small}p{0.29\textwidth}p{0.22\textwidth}p{0.39\textwidth}}
\toprule
Variable & Default/example & Meaning\\
\midrule
\endhead
OTA\_BIND\_ADDRESS & \codepath{0.0.0.0} & Listener interface; LAN access needs a non-loopback binding.\\
OTA\_PORT & \codepath{7000} & Axum HTTP port.\\
OTA\_PUBLIC\_BASE\_URL & detected or explicit & Base placed in device-facing firmware URLs. Use the PC LAN address.\\
OTA\_PROJECT\_DIR & \codepath{firmware-projects} & Safe root for extracted uploads.\\
OTA\_BUILD\_DIR & \codepath{firmware-builds} & Safe root for generated application images.\\
OTA\_DATA\_DIR & \codepath{data} & Directory containing \codepath{ota.db}.\\
OTA\_RUNTIME\_DIR & \codepath{examples/esp32-rust-ota-client} & Stable runtime template used for managed builds.\\
OTA\_MAX\_UPLOAD\_MB & \codepath{50} & Encoded project field limit in MiB.\\
OTA\_MAX\_EXTRACTED\_MB & upload limit times 10 & Maximum total declared extracted content.\\
OTA\_ALLOWED\_ORIGINS & localhost and 127.0.0.1 on port 5000 & Comma-separated browser origin allow-list.\\
OTA\_RUST\_TOOLCHAIN & \codepath{esp} & Rustup toolchain applied to builder commands; empty uses the process default.\\
WIFI\_SSID & none & Embedded Wi-Fi network name used at managed build time.\\
WIFI\_PASS & none & Embedded Wi-Fi password. Keep it in local \codepath{.env}.\\
DEVICE\_NAME & managed-device label & Friendly device name embedded into managed firmware.\\
DEVICE\_ID & MAC-derived when absent & Optional stable identity override.\\
RUST\_LOG & info defaults & Axum and tower-http tracing filter.\\
\bottomrule
\end{longtable}

\section{Command sheet}

\subsection{Start services}

\begin{lstlisting}[numbers=none]
# Windows
.\run-dev.bat

# Linux/macOS
./run-dev.sh
\end{lstlisting}

\subsection{Run OTA server checks}

\begin{lstlisting}[numbers=none]
cargo test --manifest-path ota-server/Cargo.toml
cargo clippy --manifest-path ota-server/Cargo.toml --all-targets
curl http://localhost:7000/api/ota/status
\end{lstlisting}

\subsection{Check embedded tools}

\begin{lstlisting}[numbers=none]
cargo --version
rustc +esp --version
espflash --version
rustc +esp --print target-list
\end{lstlisting}

\subsection{Build the stable runtime without flashing}

\begin{lstlisting}[numbers=none]
cd examples/esp32-rust-ota-client
cargo +esp build --release --target xtensa-esp32s3-espidf
\end{lstlisting}

\subsection{Build and flash the stable runtime}

\begin{lstlisting}[numbers=none]
cargo +esp run --release --target xtensa-esp32s3-espidf
\end{lstlisting}

The runner settings determine the exact \codepath{espflash} arguments and partition-table use. Inspect \codepath{.cargo/config.toml} in the embedded project before physical flashing.

\section{Glossary}

\begin{description}
  \item[Application image] A \codepath{.bin} containing one ESP-IDF application, suitable for an OTA app partition.
  \item[Axum] The Rust HTTP framework used by the OTA service.
  \item[Build ID] UUID identifying one compilation attempt and its durable logs.
  \item[Cargo target] The architecture and platform a Rust crate is compiled for. This project begins with \codepath{xtensa-esp32s3-espidf}.
  \item[CORS] Browser policy controlling which page origins may call the OTA API.
  \item[Deployment] A stored relationship saying one device should install one firmware record.
  \item[ELF] Linked executable format consumed by \codepath{espflash save-image}.
  \item[ESP-IDF] Espressif's official development framework and system services used beneath Rust bindings.
  \item[Firmware ID] UUID identifying one ready application image and metadata record.
  \item[Flask] The Python HTTP framework serving the dashboard and CSV API.
  \item[Heartbeat] Periodic device request refreshing version, status, IP address, and last-seen time.
  \item[Inactive slot] OTA application partition that is not currently executing and can safely receive a new image.
  \item[Managed application] Uploaded function-based Rust package composed with the stable runtime by the server.
  \item[NVS] ESP-IDF non-volatile storage, commonly used for settings and Wi-Fi configuration.
  \item[OTA] Over-the-air firmware update performed through a network connection.
  \item[Pull model] The device polls for desired firmware and initiates its own download.
  \item[Rollback] Bootloader-assisted return to a previously valid application slot when a new image fails startup validation.
  \item[SHA-256] Hash used here to confirm downloaded bytes match server metadata.
  \item[SSE] Server-sent events, the one-way HTTP stream used for CSV profiling progress.
  \item[Stable runtime] Initially flashed platform code that owns Wi-Fi, OTA, rollback, and application lifecycle.
  \item[Standalone firmware] Uploaded Cargo project built without managed-runtime composition; it owns its complete device behavior.
  \item[WAL] SQLite write-ahead logging mode, used to improve local concurrent access behavior.
  \item[ZIP slip] Archive path traversal that attempts to extract outside the intended root.
\end{description}

\section{Current implementation limits}

The first version supports ESP32-S3, local SQLite storage, polling for build logs, plain HTTP on a trusted LAN, local trusted builds, and SHA-256 integrity checks. Device serial/network terminal streaming, isolated Docker builds, multi-user authentication, publisher signing, and additional chip targets are extension work described in Chapter 12.

These limits are design boundaries, not hidden promises. Keeping them visible helps a beginner distinguish what can be used today from what still needs engineering and hardware validation.
